\documentclass[10pt,a4paper]{amsart}
\usepackage[margin=19mm]{geometry}
\usepackage{amsmath,amssymb,mathtools}
\usepackage[scr=rsfs,cal=euler]{mathalfa}
\usepackage{xfrac}
\usepackage[unicode,hidelinks,bookmarks=false]{hyperref}
\newcommand{\ie}{i.e.\ }
\newcommand{\cf}{cf.\ }
\newcommand{\resp}{respectively\ }
\newcommand{\Subsection}{\subsection}
\providecommand{\nobreakdash}{\nobreak\hbox{-}}
\setlength{\emergencystretch}{2em}
\multlinegap0pt
\usepackage[shortlabels]{enumitem}
\usepackage{bm}
\usepackage{amssymb}
\usepackage{mathtools}
\usepackage{xcolor}
\usepackage{tikz}
\newtheorem{theorem}{Theorem}
\newtheorem{proposition}{Proposition}
\newtheorem{lemma}{Lemma}
\newtheorem{claim}{Claim}
\usepackage{cleveref}
\crefname{theorem}{Theorem}{Theorems}
\crefname{proposition}{Proposition}{Propositions}
\crefname{lemma}{Lemma}{Lemmas}
\crefname{claim}{Claim}{Claims}
\newcommand\ab[1]{\lvert#1\rvert}
\newcommand\ol[1]{\overline{#1}}
\title{Color-avoiding directed paths in tournaments}
\author{Jacob Fox, Benny Sudakov,, Yuval Wigderson}
\date{Source: arXiv:2512.10438v2 (CC BY 4.0)}
\begin{document}
\maketitle

\noindent\emph{Source and licence.} Color-avoiding directed paths in tournaments. Version arXiv:2512.10438v2 (CC BY 4.0), distributed by arXiv under the Creative Commons Attribution 4.0 International licence, http://creativecommons.org/licenses/by/4.0/, as recorded in the arXiv licence metadata for this exact version and shown on the abstract page https://arxiv.org/abs/2512.10438v2. Full licence notice: \texttt{licenses/arxiv-2512.10438-CC-BY-4.0.txt}.\par\smallskip

\noindent\textit{Editorial scope.} Counted unit: the article's proposition product (source0.tex lines 562--564). The article's complete original proof of it is delivered with it (source0.tex lines 577--814, one \texttt{proof} environment of 238 lines, which the article places away from the statement). No mathematics is written, repaired or re-derived: the statement and the proof are the article's own lines, in the article's own notation, and no retained line is substituted or re-flowed.  Retained uncounted as the source-local context the proof needs: lines 168--170; lines 569--571.  Nothing was omitted from any retained span, so the package carries no omission record.  No retained line contains a citation, so the wrapper carries no bibliography and no bibliography entry.  No retained line contains a figure environment, an image-inclusion command or drawing code, so no image is delivered and the package declares no asset dependency.  Editorial formatting only, in addition: an AMS \texttt{amsart} preamble that restates the article's own declarations and macros used by the retained lines, namely the environments \texttt{theorem}, \texttt{proposition}, \texttt{lemma}, \texttt{claim} on separate counters, together with the standard packages those lines need.  The retained environments print in this document as Theorem 1, Proposition 1, Lemma 1, Claim 1 in order of appearance, whereas the article prints the same items with the numbering of its own full document.  The complete native source of the article as distributed in the arXiv e-print for this exact version is bundled unchanged as \texttt{sources/190817/source0.tex}.
\begin{theorem}\label{labelfarfromtransitive}
    If an $N$-vertex tournament is $\delta$-far from transitive, then any $q$-coloring of its edges contains a directed path of length $c\delta^2 N/q^3$ which is colored by at most three colors, where $c>0$ is an absolute constant. Furthermore, the coloring contains the square of a path such that the direction of the edges between consecutive vertices is forward, the direction of edges between vertices of distance two is backwards, and the coloring of the edges is periodic with period $3$. 
\end{theorem}

\begin{lemma}\label{lem:clean degrees}
    Let $0<\delta<\frac 12$, and let $T_0$ be an $N_0$-vertex tournament which is $\delta^2$-close to transitive. Then there exists an integer $N \geq (1-\delta)N_0$, an $N$-vertex subtournament $T$ of $T_0$, and an ordered $N$-vertex graph $G$ with the following properties. $G$ is consistent with $T$, and every vertex of $G$ has at most $4\delta N$ non-neighbors in $G$. 
\end{lemma}

\begin{proposition}\label{prop:T main product}
    We have that $\Pi_q(N) \geq c_q N^{q-Cq/\sqrt{\log q}}$, where $C>0$ is an absolute constant and $c_q>0$ is a constant depending only on $q$.
\end{proposition}

\begin{proof}[Proof of \cref{prop:T main product}]
    Recall that $\Pi_q(N) \geq 1$ for all $N$ and $q$, hence the result is vacuously true if $q$ is fixed and $C$ is large. Thus, we can and will assume henceforth that $q$ is at least some large absolute constant. The proof now proceeds by induction on $N$ for each fixed $q$; again, by picking $c_q$ sufficiently small, we can make the result vacuously true if $N$ is small, hence we have the base case of our induction. 

    We begin by defining several parameters that we will need during the proof. We let 
    \begin{equation}\label{eq:T parameters}
        p = \frac{24q}{2^{\sqrt{\log q}}}, \qquad s =2\gamma N= \frac{N}{q\cdot 2^{\sqrt{\log q}}}, \qquad \delta=\frac \gamma 8 = \frac{1}{16 q\cdot 2^{\sqrt{\log q}}}.
    \end{equation}
    To help the reader keep track of these, we note that in the proof, we will pick out $q$ disjoint sets of exactly $s$ vertices, hence we want $s$ to be smaller, but not much smaller, than $N/q$. Additionally, we will eventually encounter a special set of $p$ colors, hence we want $p$ to be smaller, but not much smaller, than $q$. The precise quantities defined above are obtained by optimizing the final contributions of the various steps of the proof. It will also be useful to notice for the future that $4\delta N = s/4$.

    We now fix some $N_0$, and assume that we have proved the result for all smaller values of $N$. We also fix a $q$-edge-colored tournament $T_0$ on $N_0$ vertices, and we aim to show that $\Pi(T_0) \geq c_q N_0^{q-Cq/\sqrt{\log q}}$. We first split into cases depending on whether or not $T_0$ is $\delta^2$-close to transitive, where we recall that $\delta$ is defined in \eqref{eq:T parameters}. If it is not, then we apply \cref{labelfarfromtransitive}. We find that $T_0$ contains a directed path of length at least $c\delta^4 N_0/q^3$, where $c>0$ is an absolute constant, which is colored by at most $3$ colors. Without loss of generality, let us assume that these three colors are $1,2,3$. This means that for every $i \in [q]\setminus\{1,2,3\}$, we have $\ell_i(T_0) \geq c\delta^4 N_0/q^3$.  As a consequence,
    \begin{align*}
        \Pi(T_0) &=\prod_{i=1}^q m_i(T) \geq \prod_{i=4}^q \min\{\ell_i(T_0),\gamma N_0\} \geq \prod_{i=4}^q \min \left\{ \frac{c\delta^4}{q^3},\gamma \right\}\cdot N_0=\left( \min \left\{ \frac{c\delta^4}{q^3},\gamma \right\} \right)^{q-3} N_0^{q-3}.
    \end{align*}
    Since $\delta$ and $\gamma$ both depend only on $q$, the first term is a positive constant depending only on $q$, and in particular can be made larger than $c_q$ by choosing $c_q$ appropriately. Moreover, by picking $C$ appropriately, we can ensure that $Cq/\sqrt{\log q}\geq 3$ for all $q$, hence $N_0^{q-3}\geq N_0^{q-Cq/\sqrt{\log q}}$. Therefore, we conclude that $\Pi(T_0) \geq c_q N_0^{q-Cq/\sqrt{\log q}}$, as desired. 

    We may thus assume that $T_0$ is $\delta^2$-close to transitive. We now apply \cref{lem:clean degrees} to pass to a subtournament $T$ on $N \geq (1-\delta)N_0$ vertices, as well as an ordered graph $G$ consistent with $T$ in which every vertex has at most $4\delta N$ non-neighbors. As $T$ is a subtournament of $T_0$, we have $\Pi(T_0)\geq \Pi(T)$, and therefore it suffices to prove that $\Pi(T) \geq c_q N_0^{q-Cq/\sqrt{\log q}}$. Moreover, since $N \geq (1-\delta)N_0\geq 2^{-2\delta}N_0$, it in turn suffices to prove that 
    \begin{equation}\label{eq:goal}
    \Pi(T)\geq 2^{\delta q}\cdot c_q\cdot N^{q-Cq/\sqrt{\log q}},
    \end{equation}
    which will be our goal for the remainder of the proof. We recall that our inductive hypothesis is that for all $N'<N$, we have $\Pi_q(N') \geq c_q\cdot (N')^{q-Cq/\sqrt{\log q}}$. 


    We say that a color $i\in [q]$ is \emph{long} if $\ell_i(T)\geq \gamma N$, and we say that it is \emph{short} otherwise. Note that for any color $i \in [q]$, we have 
    \[
        m_i(T) = \min\{\ell_i(T), \gamma N\} = 
        \begin{cases}
            \ell_i(T) & \text{if $i$ is short},\\
            \gamma N & \text{if $i$ is long}.
        \end{cases}
    \]
    The long colors play a special role in our argument, and must be dealt with separately. In most of what follows, we argue almost exclusively about the short colors. 

    As $G$ is an $N$-vertex ordered graph, we henceforth identify the vertices of $G$ with the interval $[N]$, so that the ordering $\prec$ of $G$ agrees with the standard ordering of $[N]$. We alternate between these two notations as convenient.
    We begin by partitioning $[N]=V(G)=V(T)$ into four intervals $A,B,C,D$, of lengths $\frac N2-4s, 4s, 4s, \frac N2-4s$, respectively, where we recall that $s$ was defined in \eqref{eq:T parameters}. That is, we have
    \[
        A = [1,\tfrac N2-4s], \qquad B = (4s,\tfrac N2], \qquad C = (\tfrac N2, \tfrac N2+4s], \qquad D = (\tfrac N2+4s,N].
    \]
    For a vertex $v \in A\cup B$ and a color $i\in [q]$, we denote by $\ell_i(^\to v)$ the length of the longest directed path in $T[A \cup B]$ which ends at the vertex $v$ and whose edges avoid the color $i$. Note that we are restricting only to those directed paths which are entirely contained in $A\cup B$. Similarly, for $v\in C\cup D$, we denote by $\ell_i(v^\to)$ the length of the longest $i$-avoiding directed path in $T[C\cup D]$ that starts at the vertex $v$.

    For every color $i \in [q]$, we let $X_i$ consist of the $s$ vertices in $A \cup B$ which have the largest values of $\ell_i({}^\to v)$. That is, for every vertex $v \in A \cup B$, we compute the integer $\ell_i(^\to v)$, we rank the vertices of $A \cup B$ according to these integers, and then let $X_i$ comprise the $s$ top-most vertices in this ranking. If there are ties in the ranking, we break them arbitrarily. Similarly, we define $Y_i$ to comprise the $s$ vertices in $C \cup D$ with the largest values of $\ell_i(v^\to)$. By definition, $\ab{X_i}=\ab{Y_i}=s$ for all $i$. 

    

    Let $0\leq r\leq q$ be the number of long colors in $T$. We split the $q-r$ short colors into several further types, as follows.
    Let $i$ be a short color. We call the color $i$ \emph{left-condensed} if $\ab{X_i \cap B} \geq s/2$, and \emph{left-diffuse} if $\ab{X_i\cap B}<s/2$. Similarly, $i$ is \emph{right-condensed} if $\ab{Y_i\cap C}\geq s/2$, and \emph{right-diffuse} otherwise. Finally, we call color $i$ \emph{condensed} if it is both left-condensed and right-condensed. We stress that we only apply these terms to short colors; in particular, whenever we speak of a condensed or left-diffuse color in what follows, that color will always also be short. 
    We now split into cases depending on the number of condensed colors, recalling that there are exactly $r$ long colors and $q-r$ short colors. 

    \paragraph{\bf Case 1: There are at least $\bm{q-r-2p}$ condensed colors} 
    Let $\Lambda \subseteq [q]$ denote the set of long colors, and let $\Xi \subseteq [q] \setminus \Lambda$ denote the set of condensed colors.
    Recall that every vertex of $G$ has at most $4\delta N = s/4$ non-neighbors. As a consequence, for every $i \in \Xi$, the number of edges of $G$ between $X_i \cap B$ and $Y_i \cap C$ is at least
    \[
        \sum_{v \in X_i \cap B} \left(\ab{Y_i \cap C}-\frac s4\right) \geq \sum_{v \in X_i\cap B} \left( \frac s2-\frac s4 \right) =\frac s4 \ab{X_i \cap B} \geq \frac{s^2}{8},
    \]
    where we use that color $i$ is right-condensed in the first inequality, and that it is left-condensed in the final inequality.
    
    Let us call a color $i \in \Xi$ \emph{useless} if all edges of $G$ between $X_i \cap B$ and $Y_i \cap C$ are of color $i$, and \emph{useful} otherwise. By the previous computation, if $i$ is useless, then there are at least $s^2/8$ color-$i$ edges between $B$ and $C$. Since there are at most $\ab B \ab C = 16s^2$ edges of $G$ between $B$ and $C$, we conclude that the number of useless colors is at most $(16s^2)/(s^2/8)=128$. 

    By definition, for every useful color $i$, there is at least one edge $v_iw_i\in E(G)$, with $v_i \in X_i \cap B$ and $w_i \in Y_i \cap C$, which does not receive color $i$ (for if there were no such edge, then $i$ would be useless). As $G$ is consistent with $T$ and $B$ precedes $C$ in the ordering, we know that the edge $v_iw_i$ is directed as $v_i\to w_i$. These facts imply that $T$ contains an $i$-avoiding directed path of length $\ell_i(^\to v_i)+\ell_i(w_i^\to)$, obtained by taking the longest $i$-avoiding directed path in $T[A\cup B]$ ending at $v_i$, concatenating it with the edge $v_i\to w_i$, and then concatenating that with the longest $i$-avoiding directed path in $T[C\cup D]$ starting at $w_i$. This is indeed a directed path since $v_i \to w_i$ and since the two paths we are concatenating have disjoint vertex sets; additionally, since the edge $v_i w_i$ does not receive color $i$, this is indeed an $i$-avoiding directed path.

    Let $T'$ be the subtournament of $T$ induced on the vertex set $(A \cup B) \setminus \bigcup_{i=1}^q X_i$. Note that $T'$ has at least $N'=\frac N2 - sq$ vertices. Moreover, for the special vertex $v_i$ found above, we claim that $\ell_i(^\to v_i) \geq \ell_i(T')$. Indeed, the longest $i$-avoiding directed path in $T'$ is entirely contained in $V(T')\subseteq A\cup B$, and ends at some vertex $u \in V(T')$, which in particular does not belong to $X_i$, as $X_i$ is disjoint from $V(T')$. As $X_i$ comprises the $s$ vertices with the highest values of $\ell_i(^\to \bullet)$, we see that $\ell_i(^\to v_i) \geq \ell_i(^\to u)=\ell_i(T')$.  Similarly, if we let $T''$ be the subtournament induced on $(C \cup D) \setminus \bigcup_{i=1}^q Y_i$, then $\ell_i(w_i^\to) \geq \ell_i(T'')$ for each $i \in [q]$. 

    Now, for each useful color $i \in \Xi$, we have that
    \[
        \ell_i(T) \geq \ell_i(^\to v_i)+\ell_i(w_i^\to) \geq \ell_i(T')+\ell_i(T'') \geq 2\sqrt{\ell_i(T') \ell_i(T'')},
    \]
    where the final bound is the inequality of arithmetic and geometric means. Since every such color is short, we also know that $m_i(T)=\min\{\ell_i(T),\gamma N\}=\ell_i(T)$, hence
    \[
        m_i(T)=\ell_i(T)\geq 2 \sqrt{\ell_i(T')\ell_i(T'')}\geq 2\sqrt{m_i(T')m_i(T'')}
    \]
    for every useful $i \in \Xi$.
    Similarly, for every long color $i \in \Lambda$, we have that $m_i(T)=\min\{\ell_i(T),\gamma N\}=\gamma N$, hence
    \[
        m_i(T) = \gamma N \geq \gamma \ab{T'}+\gamma \ab{T''} \geq m_i(T')+m_i(T'') \geq 2\sqrt{m_i(T')m_i(T'')}.
    \]
    In other words, we have found that for at least $q-2p-128$ of the colors $i$ (namely the useful or long colors), the inequality $m_i(T) \geq 2 \sqrt{m_i(T')m_i(T'')}$ holds. The remaining $2p+128$ colors are short, hence satisfy $m_i(T)=\ell_i(T)$. Moreover, for these remaining colors, we trivially have the inequality
    \[
        m_i(T)=\ell_i(T) \geq \max\{\ell_i(T'),\ell_i(T'')\} \geq \sqrt{\ell_i(T')\ell_i(T'')}\geq \sqrt{m_i(T')m_i(T'')}.
    \]
    Putting this all together, we find that
    \[
        \Pi(T) = \prod_{i=1}^q m_i(T) \geq 2^{q-2p-128} \prod_{i=1}^q\sqrt{m_i(T')m_i(T'')} = 2^{q-2p-128} \sqrt{\Pi(T')\Pi(T'')}.
    \]
    Finally, recall that $T'$ and $T''$ are both $q$-edge-colored tournaments on at least $N'=\frac N2-sq$ vertices. By the definition of $\Pi_q(N')$, this shows that $\Pi(T'), \Pi(T'') \geq \Pi_q(N')$. In conclusion, we have that
    \[
        \Pi(T) \geq 2^{q-2p-128} \Pi_q(N').
    \]
    By the inductive hypothesis, we know that $\Pi_q(N')\geq c_q (N')^{q-Cq/\sqrt{\log q}}$.
    Additionally, by the definition of $s$ in \eqref{eq:T parameters}, we know that $N' = (\frac 12-\frac{1}{2^{\sqrt{\log q}}})N$. Since we may assume that $q$ is sufficiently large, we have that $2^{\sqrt{\log q}}\geq \sqrt{\log q}$, and thus $N'\geq (\frac 12-\frac{1}{{\sqrt{\log q}}})N$.  Therefore,
    \begin{align*}
        \frac{(N')^{q-Cq/\sqrt{\log q}}}{N^{q-Cq/\sqrt{\log q}}} &\geq \left( \frac 12-\frac{1}{{\sqrt{\log q}}} \right)^{q-Cq/\sqrt{\log q}} = 2^{-q+Cq/\sqrt{\log q}} \left( 1-\frac{2}{{\sqrt{\log q}}} \right)^{q-Cq/\sqrt{\log q}}.
    \end{align*}
    Using the inequality $1-x \geq 2^{-2x}$, valid for all $x \in [0,\frac 12]$ (and recalling that $q$ is at least a sufficiently large constant, so that we may apply this inequality), we have that
    \begin{align*}
        \left( 1-\frac{2}{{\sqrt{\log q}}} \right)^{q-Cq/\sqrt{\log q}} \geq\left( 1-\frac{2}{{\sqrt{\log q}}} \right)^q \geq 2^{-4q/\sqrt{\log q}}.
    \end{align*}
    Combining these bounds, we conclude that
    \begin{align*}
        \frac{\Pi(T)}{c_q N^{q-Cq/\sqrt{\log q}}} &\geq \frac{2^{q-2p-128} \Pi_q(N')}{c_q N^{q-Cq/\sqrt{\log q}}}\geq 2^{q-2p-128} \cdot \frac{(N')^{q-Cq/\sqrt{\log q}}}{N^{q-Cq/\sqrt{\log q}}}\\
        &\geq 2^{q-2p-128} \cdot 2^{-q+Cq/\sqrt{\log q}} \cdot 2^{-4q/\sqrt{\log q}}\\
        &= 2^{(C-4)q/\sqrt{\log q} -2p -128}.
    \end{align*}
    Finally, we recall from \eqref{eq:T parameters} that $p = 24q/2^{\sqrt{\log q}} \leq q/\sqrt{\log q}$, where the final inequality holds since $q$ is sufficiently large. Therefore, the exponent above is at least $(C-6)q/\sqrt{\log q}-128$. In particular, as $C$ is sufficiently large, we have that this exponent is at least $2q/{\sqrt{\log q}}\geq 2\delta q$ for all $q$. We conclude that $\Pi(T) \geq 2^{2\delta q}\cdot c_q\cdot N^{q-Cq/\sqrt{\log q}}$. 
    This is precisely the inequality \eqref{eq:goal} that we set out to prove, and thus this completes the proof in Case 1.

    \paragraph{\bf Case 2: There are fewer than $\bm{q-r-2p}$ condensed colors.} In this case, there are at least $2p$ non-condensed colors, implying that there are at least $p$ left-diffuse colors or at least $p$ right-diffuse colors. The two cases are handled identically (in fact they are equivalent upon reversing the orientation of $T$), so we may assume without loss of generality that there are at least $p$ left-diffuse colors. We let $\Delta_L\subseteq [q]$ be a set of exactly $p$ left-diffuse colors.

    The proof in this case relies repeatedly on the following simple, but remarkably useful, observation, which follows immediately from the definition of the set $X_i$. 
    \begin{claim}\label{claim:exceptional}
        Let $i$ be a short color, and let $v \in X_i$. There is a set $E_i(v) \subseteq A \cup B$ with $\ab{E_i(v)}< 2s$ such that the following holds for all $w \in (A \cup B) \setminus  E_i(v)$. The vertices $v$ and $w$ are adjacent in $G$, and if $v\to w$ in $T$, then the color of the edge $vw$ is $i$.
    \end{claim}
    Here, we think of the set $E_i(v)$ as a set of ``exceptional'' vertices: apart from this small set of vertices, every out-directed edge of $v$ to $A \cup B$ receives color $i$. 
    \begin{proof}
        Let $\ol N(v)$ be the non-neighborhood of $v$ in $G$.
        We fix an $i$-avoiding directed path $P$ in $T[A\cup B]$ ending at $v$ of length $\ell_i(^\to v)$. As the color $i$ is short, we must have $\ab{V(P)}=\ell_i(^\to v)<\gamma N$. Now let $E_i(v)=V(P) \cup \ol N(v) \cup X_i$; we claim that this definition satisfies the desired property. 

        First, we have that $\ab{E_i(v)}\leq \ab{V(P)}+\ab{\ol N(v)}+\ab{X_i} \leq \gamma N+4\delta N+s < 2s$. Additionally, as $\ol N(v)\subseteq E_i(v)$, we certainly have that $v$ is adjacent in $G$ to all vertices in $(A\cup B) \setminus E_i(v)$. For the final property, fix some $w \in (A\cup B) \setminus E_i(v)$ with $v\to w$. As $w \notin X_i$ (since $X_i \subseteq E_i(v)$), we have that
        $\ell_i(^\to v) \geq \ell_i(^\to w)$, for $X_i$ consists of the $s$ vertices with the largest value of $\ell_i(^\to \bullet)$. Moreover, as $w \notin E_i(v)$, we see that $w \notin V(P)$. Consider the path $P +w$, obtained from $P$ by adding the edge $v\to w$ as the final edge. This is indeed a directed path, since $v\to w$ and $w \notin V(P)$. Additionally, if the edge $vw$ does not receive color $i$, then $P+w$ is an $i$-avoiding directed path in $T[A \cup B]$ ending at $w$. Therefore its length is at most $\ell_i(^\to w)$, from which we find
        \[
            \ell_i(^\to v) \geq \ell_i(^\to w) \geq \ab{V(P+w)} = \ab{V(P)}+1 = \ell_i(^\to v)+1,
        \]
        a contradiction. Thus, the edge $vw$ must receive color $i$. 
    \end{proof}

    For $i \in \Delta_L$, let $U_i = X_i \cap A$. As the color $i$ is left-diffuse, we know that $\ab{X_i \cap B}<s/2$, hence $\ab{U_i}\geq s/2$. 
    The following immediate consequence of \cref{claim:exceptional} shows that the sets $\{U_i : i \in \Delta_L\}$ are pairwise disjoint.
    \begin{claim}\label{claim:T Ui disjoint}
        Let $i,j \in \Delta_L$ be left-diffuse colors with $i \neq j$. Then the sets $U_i$ and $U_j$ are disjoint.
    \end{claim}
    \begin{proof}
        Suppose for contradiction that there is some vertex $v \in U_i \cap U_j$. Recall that $\ab B=4s$, and note that
        \[
            \ab{E_i(v)} + \ab{E_j(v)} < 2s+2s = 4s = \ab B,
        \]
        by \cref{claim:exceptional}.
        This implies that there is some vertex $w \in B \setminus (E_i(v) \cup E_j(v))$. By \cref{claim:exceptional}, we have $vw \in E(G)$. Moreover, since $v \in U_i \subseteq A$ and $w \in B$, we know that $v \prec w$, hence must have that $v\to w$ in $T$. Thus, by \cref{claim:exceptional}, the edge $vw$ receives color $i$. 
        But for the exact same reason, it must receive color $j$, a contradiction.
    \end{proof}
    Let $U = \bigcup_{i \in \Delta_L}U_i$. 
    By definition, for every $v \in U$, we must have that $v \in U_i$ for some $i\in \Delta_L$; moreover, this $i$ is unique by \cref{claim:T Ui disjoint}. Thus, for a vertex $v \in U$, we denote by $\iota(v)$ the unique index $i\in \Delta_L$ for which $v \in U_{i}$.
    
    Thanks to \cref{claim:exceptional}, we know a great deal about the colors of edges within $U$: almost all such edges are colored according to the index $\iota$ of their left endpoint. In particular, almost all edges in $U$ are colored with colors from $\Delta_L$. 
    In most of the remainder of the proof, we will use this structure to show that for all colors $k \notin \Delta_L$, we can find a long $k$-avoiding directed path within $U$; this is plausible, since we know that the edges within $U$ that \emph{are} colored $k$ must be quite restricted.

    Our first goal in this direction is to construct an appropriate structure within $U$ that we will use to glue long $k$-avoiding paths together. The precise structure we use is given in the following claim. For a vertex $v$ and vertex subset $S$, we write $v \prec S$ if $v \prec w$ for all $w \in S$, and we write $S \prec v$ if $w \prec v$ for all $w \in S$. 
    \begin{claim}\label{claim:gluing structure}
        There exist vertices $v_1,\dots,v_{t+1} \in U$, as well as sets $S_1,\dots,S_{t}\subseteq U$, with the following properties.
        \begin{enumerate}[label=(\roman*)]
            \item We have $t=p/24$ and $\ab{S_a}\geq s$ for all $a \in [t]$.\label{it:sizes}
            \item We have $v_1 \prec S_1 \prec v_2 \prec S_2 \prec v_3 \prec \dots \prec v_{t} \prec S_{t}\prec v_{t+1}$.\label{it:order}
            \item $S_a$ is contained in the out-neighborhood of $v_a$ and in the in-neighborhood of $v_{a+1}$, for all $a \in [t]$. That is, the edges are directed as $v_a \to S_a \to v_{a+1}$. \label{it:direction}
            \item All edges from $v_a$ to $S_a$ and all edges from $S_a$ to $v_{a+1}$ are colored with a color from $\Delta_L$.\label{it:color}
        \end{enumerate}
    \end{claim}
    \begin{proof}
        We begin by setting $v_1$ to be the first vertex in $U$ according to $\prec$. We now recursively define $S_a$ and $v_{a+1}$, given the value of $v_a$, such that the desired properties hold. Having already defined $v_1$, we can start the recursion.

        So suppose the value of $v_a$ is given. If there are fewer than $12s$ vertices in $U$ which come after $v_a$ under $\prec$, then we stop the recursion. Otherwise, we let $I_a$ denote the next $4s$ vertices in $U$ after $v_a$, and let $J_a$ be the subsequent $8s$ vertices in $U$ after $I_a$. We finally define $I_a' = I_a \setminus E_{\iota(v_a)}(v_a)$. 
        Note that 
        \[
            4s = \ab{I_a} \geq \ab{I_a'} = \ab{I_a \setminus E_{\iota(v_a)}(v_a)} \geq \ab{I_a} - \ab{E_{\iota(v_a)}(v_a)} \geq 4s - 2s = 2s,
        \]
        that is, that $2s \leq \ab{I_a'} \leq 4s$.

        Next, consider an auxiliary bipartite graph $\Gamma$ between $I_a'$ and $J_a$, where we join $v \in I_a'$ to $w \in J_a$ if and only if $w \in E_{\iota(v)}(v)$. As $\ab{E_{\iota(v)}(v)}\leq 2s$ by \cref{claim:exceptional}, we see that in $\Gamma$, all vertices in $I_a'$ have degree at most $2s$. Thus, the number of edges in $\Gamma$ is at most $2s\ab{I_a'} \leq 8s^2$. As a consequence, there must exist some vertex $w \in J_a$ whose degree in $\Gamma$ is at most $(8s^2)/\ab{J_a}=s$. We let $v_{a+1}$ be such a choice of $w \in J_a$. Finally, we let $S_a\subseteq I_a'$ be the non-neighbors of $v_{a+1}$ in $\Gamma$. By construction, we have
        \[
            \ab{S_a} \geq \ab{I_a'}-s \geq 2s-s=s.
        \]
        We have thus finished defining $S_a$ and $v_{a+1}$, and it remains to verify the claimed properties. We have just shown that $\ab{S_a}\geq s$, as claimed in \ref{it:sizes}. As $v_a \prec I_a \prec J_a$, we certainly have $v_a \prec S_a \prec v_{a+1}$, as claimed in \ref{it:order}. As $S_a \subseteq I_a' \subseteq (A\cup B) \setminus E_{\iota(v_a)}(v_a)$, we conclude from \cref{claim:exceptional} that all edges from $v_a$ to $S_a$ are oriented as $v_a \to S_a$, and they all receive color $\iota(v_a)$, proving the first halves of \ref{it:direction} and \ref{it:color}. For the corresponding claims about edges from $S_a$ to $v_{a+1}$, fix some $v \in S_a$. By the choice of $v_{a+1}$, we know that $v_{a+1} \notin E_{\iota(v)}(v)$, hence we again conclude that the edge $vv_{a+1}$ is oriented as $v \to v_{a+1}$, and that it receives color $\iota(v)\in \Delta_L$, completing the proofs of \ref{it:direction} and \ref{it:color}.

        Note that at every step of this process, we remove at most $12s$ vertices from $U$, hence we can continue this so long as $a<\ab U/(12s)$. Also note that $U$ is the union of the $p$ disjoint sets $U_i$ with $i \in \Delta_L$, and each such set has size at least $s/2$, by the definition of a left-diffuse color. This implies that $\ab U \geq ps/2$, and hence we can continue the process at least to step $t$, where
        \[
            t = \frac{\ab U}{12 s} \geq \frac{ps}{24s} = \frac{p}{24}.\qedhere
        \]
    \end{proof}



    For each $a \in [t]$, let $T_a$ be the subtournament of $T$ induced on the vertex set $S_a$. The key claim to complete the proof is the following, which allows us to glue together $k$-avoiding paths in all of the tournaments $T_a$, for every color $k \notin \Delta_L$. 
    \begin{claim}\label{claim:T glue paths}
        Let $k\in [q] \setminus \Delta_L$ be a color. We have that
        \[
            m_k(T) \geq \sum_{a=1}^t m_k(T_a).
        \]
    \end{claim}
    \begin{proof}
        Let $P_a$ be a longest $k$-avoiding directed path in $T_a$, for all $a \in [t]$. We concatenate these paths into a longer path $P$ defined as $v_1P_1v_2P_2 \dots P_t v_{t+1}$. That is, we start at $v_1$, go to the first vertex of the directed path $P_1$, traverse $P_1$, go to $v_2$, and then continue in this fashion. By \cref{claim:gluing structure}, this gives us a directed path, since $v_a \to S_a \to v_{a+1}$ for all $a \in [t]$. Moreover, all edges $v_a \to S_a \to v_{a+1}$ are colored by colors in $\Delta_L$, by \cref{claim:gluing structure}, and in particular are not colored by the color $k$. Thus, the path $P$ avoids the color $k$, and its length is at least $\sum_{a=1}^t \ell_k(T_a)$; this proves that
        \[
            \ell_k(T) \geq \sum_{a=1}^t \ell_k(T_a).
        \]
        We now split into cases depending on whether $k$ is long or short. If $k$ is short, then $m_k(T)=\ell_k(T)$, hence
        \[
            m_k(T)=\ell_k(T) \geq \sum_{a=1}^t \ell_k(T_a) \geq \sum_{a=1}^t m_k(T_a).
        \]
        On the other hand, if $k$ is long, then
        \[
            m_k(T) = \gamma N \geq \sum_{a=1}^t \gamma \ab{T_a} \geq \sum_{a=1}^t m_k(T_a),
        \]
        where the first inequality uses that the sets $S_1,\dots,S_t$ are pairwise disjoint subsets of $[N]$. 
        In either case, we have the claimed bound.
    \end{proof}
    From \cref{claim:T glue paths} and the inequality of arithmetic and geometric means, we find that
    \[
        m_k(T) \geq \sum_{a=1}^t m_k(T_a) \geq t \left( \prod_{a=1}^t m_k(T_a) \right)^{1/t}
    \]
    for every color $k \in [q] \setminus \Delta_L$. On the other hand, for every color $k \in \Delta_L$, we trivially have
    \[
        m_k(T) \geq \max\{m_k(T_1),\dots,m_k(T_t)\} \geq \left( \prod_{a=1}^t m_k(T_a) \right)^{1/t}.
    \]
    Combining these bounds, we find that
    \begin{align*}
        \Pi(T) &= \prod_{k=1}^q m_k(T) =\prod_{k\in \Delta_L} m_k(T) \cdot \prod_{k\in [q] \setminus \Delta_L} m_k(T)\\
        &\geq \prod_{k\in \Delta_L} \left( \prod_{a=1}^t m_k(T_a) \right)^{1/t} \cdot \prod_{k\in [q]\setminus \Delta_L} t \left( \prod_{a=1}^t m_k(T_a) \right)^{1/t}\\
        &= t^{q-p} \left( \prod_{a=1}^t \prod_{k=1}^q m_k(T_a) \right)^{1/t}\\
        &=t^{q-p} \left( \prod_{a=1}^t \Pi(T_a) \right)^{1/t}\\
        &\geq t^{q-p}\cdot \Pi_q(s),
    \end{align*}
    where the final step uses that each $T_a$ is a $q$-edge-colored tournament on $\ab{S_a}\geq s$ vertices.

    By the induction hypothesis on $s$ vertices, we know that $\Pi_q(s) \geq c_q s^{q-Cq/\sqrt{\log q}}$. From the definition of $s$ in \eqref{eq:T parameters}, we have that
    \[
        \frac{s^{q-Cq/\sqrt{\log q}}}{N^{q-Cq/\sqrt{\log q}}} = \left( \frac{1}{q\cdot 2^{\sqrt{\log q}}} \right)^{q-Cq/\sqrt{\log q}} = q^{-q+Cq/\sqrt{\log q}} \cdot 2^{-q\sqrt{\log q}+Cq} \geq q^{-q} \cdot 2^{(C-1)q\sqrt{\log q}}.
    \]
    Similarly, since $t=p/24$ from \cref{claim:gluing structure}, and from the definition of $p$ from \eqref{eq:T parameters}, we have that
    \[
        t^{q-p} = \left( \frac{q}{2^{\sqrt{\log q}}} \right)^{q-p} =q^{q-p} \cdot 2^{-q \sqrt{\log q}+p\sqrt{\log q}} \geq q^{q-p} \cdot 2^{-q\sqrt{\log q}}.
    \]
    Since $q$ is sufficiently large, we have that $2^{\sqrt{\log q}}\geq 24\sqrt{\log q}$, and hence $p \leq q/\sqrt{\log q}$. Therefore,
    \[
        t^{q-p}\geq q^{q-p} \cdot 2^{-q\sqrt{\log q}} \geq q^{q-q/\sqrt{\log q}} \cdot 2^{-q\sqrt{\log q}}=q^q \cdot 2^{-2q\sqrt{\log q}}. 
    \]
    Combining these bounds, we find that
    \begin{align*}
        \frac{\Pi(T)}{c_q N^{q-Cq/\sqrt{\log q}}} &\geq \frac{t^{q-p}\cdot \Pi_q(s)}{c_q N^{q-Cq/\sqrt{\log q}}} \geq t^{q-p}\cdot \frac{s^{q-Cq/\sqrt{\log q}}}{N^{q-Cq/\sqrt{\log q}}}\\
        &\geq \left( q^q \cdot 2^{-2q \sqrt{\log q}} \right) \left( q^{-q} \cdot 2^{(C-1)q\sqrt{\log q}} \right)\\
        &= 2^{(C-3)q\sqrt{\log q}}. 
    \end{align*}
    In particular, as $C$ is sufficiently large, then this exponent is at least $2\delta q$, hence we find that $\Pi(T) \geq 2^{\delta q}\cdot c_q\cdot N^{q-Cq/\sqrt{\log q}}$, as we wanted to prove in \eqref{eq:goal}. This completes the proof of Case 2, and thus of \cref{prop:T main product}.
\end{proof}

\end{document}
